\documentclass[11pt,a4paper]{article}

% ---- Packages ----
\usepackage[utf8]{inputenc}
\usepackage[T1]{fontenc}
\usepackage{amsmath,amssymb,amsthm}
\usepackage{enumitem}
\usepackage{geometry}
\usepackage{booktabs}
\usepackage{xcolor}
\usepackage{titlesec}
\usepackage{hyperref}
\usepackage{parskip}
\usepackage{mdframed}

\geometry{margin=2.2cm}

% ---- Colours ----
\definecolor{blockcolor}{RGB}{0,70,130}
\definecolor{hintcolor}{RGB}{140,90,0}
\definecolor{exbg}{RGB}{245,245,235}
\definecolor{anscolor}{RGB}{200,200,200}

% ---- Book example environment ----
\newmdenv[
  backgroundcolor=exbg,
  linecolor=blockcolor!40,
  linewidth=1pt,
  roundcorner=4pt,
  innertopmargin=10pt,
  innerbottommargin=10pt,
  innerleftmargin=12pt,
  innerrightmargin=12pt,
  skipabove=8pt,
  skipbelow=8pt
]{bookexample}

% ---- Answer box environment ----
\newmdenv[
  linecolor=anscolor,
  linewidth=0.8pt,
  roundcorner=3pt,
  innertopmargin=8pt,
  innerbottommargin=8pt,
  innerleftmargin=10pt,
  innerrightmargin=10pt,
  skipabove=6pt,
  skipbelow=6pt,
  backgroundcolor=white
]{answerbox}

% ---- Section formatting ----
\titleformat{\section}{\large\bfseries\color{blockcolor}}{}{0em}{}[\vspace{-0.3em}\rule{\textwidth}{0.4pt}]

% ---- Header ----
\title{\textbf{Tutorial 5 --- Potential Outcomes, Causality,\\ and Counterfactual Reasoning}\\\medskip \Large Questions}
\author{Introductory Econometrics --- TA Session}
\date{Duration: 50 minutes \\ \medskip \small Based on Wooldridge, Sections 1-4 and 2-7a}

\begin{document}
\maketitle
\thispagestyle{empty}

%=============================================================================
\section{Block 1: Motivation --- Why Do We Need Potential Outcomes? \textnormal{(5 min)}}
%=============================================================================

\begin{bookexample}
\textbf{\color{blockcolor}Causality, Ceteris Paribus, and Counterfactual Reasoning} {\small (Section 1-4)}

\smallskip
In economics we want to know whether one variable has a \textbf{causal effect} on another. This requires \textbf{ceteris paribus} reasoning: what happens to an outcome when we change one factor, \emph{holding all other factors fixed}. Since we can rarely hold other factors fixed, we resort to \textbf{counterfactual reasoning}: ``what \emph{would have happened} under a different state of the world?'' The challenge is that we can never observe the same unit in both states simultaneously. This is the \textbf{fundamental problem of causal inference}. The textbook illustrates this with four examples:

\smallskip
\begin{itemize}[leftmargin=1.5em,itemsep=3pt]
  \item \textbf{Ex.~1.3 --- Fertilizer on Crop Yield.} A farmer obtains 180 bushels/acre using fertilizer. The counterfactual (yield \emph{without} fertilizer on the same plot, same season) is unobservable. An ideal experiment randomly assigns fertilizer to identical plots. In practice, farmers choose fertilizer based on soil quality --- a \textbf{confounding factor} that also affects yield.

  \item \textbf{Ex.~1.4 --- Return to Education.} If a person gets one more year of education, by how much does their wage rise? A social planner could randomly assign education levels, but this is infeasible. People \emph{choose} education based on ability and background, so comparing wages across education levels confounds the education effect with pre-existing differences (\textbf{selection bias}).

  \item \textbf{Ex.~1.5 --- Police and Crime.} Does hiring more police reduce crime? Cities with high crime already hire more police, creating a positive correlation even if police \emph{reduce} crime. This is \textbf{reverse causality} (simultaneity).

  \item \textbf{Ex.~1.6 --- Minimum Wage and Unemployment.} Political and economic forces that set the minimum wage also affect employment. It is impossible to isolate the causal effect without holding these \textbf{confounding forces} fixed.
\end{itemize}

\smallskip
In each case, the core problem is the same: we observe one state of the world but need to compare it with a counterfactual we cannot see. The \textbf{potential outcomes framework} (Section~2-7a) formalizes this problem mathematically, and \textbf{random assignment} provides the cleanest solution.
\end{bookexample}

\textbf{Question 1} (Quick warm-up)

\begin{enumerate}[label=(\alph*)]
\item In Example~1.4, why can we \emph{not} simply compare the average wage of people with 16~years of education to that of people with 15~years and call the difference the causal return to education?

\begin{answerbox}
\vspace{3cm}
\end{answerbox}

\item For Examples~1.3--1.6, each describes an ideal experiment that is infeasible. What is the \emph{one} feature all these ideal experiments share that would solve the causal inference problem?

\begin{answerbox}
\vspace{3cm}
\end{answerbox}
\end{enumerate}

\newpage
%=============================================================================
\section{Block 2: Potential Outcomes Notation \textnormal{(13 min)}}
%=============================================================================

\begin{bookexample}
\textbf{\color{blockcolor}Section 2-7a --- Counterfactual (or Potential) Outcomes, Causality, and Policy Analysis}

\smallskip
Consider a binary treatment: $x_i = 1$ if unit $i$ is in the treatment group, $x_i = 0$ if in the control group. For each unit~$i$ there are two \textbf{potential outcomes}: $y_i(1)$ (outcome if treated) and $y_i(0)$ (outcome if not treated). Both exist conceptually, but we only observe one.

The \textbf{individual treatment effect} is $\tau_i = y_i(1) - y_i(0)$.

The \textbf{Average Treatment Effect (ATE)} and \textbf{Average Treatment Effect on the Treated (ATT)} are:
\begin{equation}
  ATE = E[y(1)] - E[y(0)] \tag{2.75}
\end{equation}
\begin{equation}
  ATT = E[y(1) - y(0) \mid x = 1] \tag{2.76}
\end{equation}

The \textbf{observed outcome} is given by the ``switching equation'':
\begin{equation}
  y_i = x_i \cdot y_i(1) + (1 - x_i) \cdot y_i(0) \tag{2.77}
\end{equation}

Given a random sample, we observe only one of $y_i(0)$ and $y_i(1)$: the treatment ``switches'' which potential outcome we see.
\end{bookexample}

\bigskip
\textbf{Question 2} (Notation --- Definitions)

\medskip
Consider a binary treatment $x_i \in \{0,1\}$ for individual~$i$.

\begin{enumerate}[label=(\alph*)]
\item Define the two \textbf{potential outcomes} $y_i(0)$ and $y_i(1)$. What do they represent?

\begin{answerbox}
\vspace{3.5cm}
\end{answerbox}

\item Define the \textbf{individual treatment effect} $\tau_i$. Why can we never compute it directly?

\begin{answerbox}
\vspace{3.5cm}
\end{answerbox}

\item Write the \textbf{observed outcome} $y_i$ in terms of $x_i$, $y_i(0)$, and $y_i(1)$ (equation~[2.77]). Verify the formula for each value of $x_i$.

\begin{answerbox}
\vspace{4cm}
\end{answerbox}
\end{enumerate}

\newpage
%-----------------------------------------------------------------------------
\textbf{Question 3} (Math --- Algebraic manipulation)

\medskip
Starting from $y_i = x_i \cdot y_i(1) + (1 - x_i)\cdot y_i(0)$:

\begin{enumerate}[label=(\alph*)]
\item Show that this is equivalent to:
\[
  y_i = y_i(0) + \bigl[y_i(1) - y_i(0)\bigr]\cdot x_i = y_i(0) + \tau_i \cdot x_i
\]

\begin{answerbox}
\vspace{5cm}
\end{answerbox}

\item Now suppose the treatment effect is \textbf{constant}: $\tau_i = \tau$ for all~$i$. Let $\beta_0 = E[y(0)]$ and $u_i = y_i(0) - E[y(0)]$. Show that:
\[
  y_i = \beta_0 + \tau \cdot x_i + u_i
\]
What does this equation look like? What is the connection to regression?

\begin{answerbox}
\vspace{5cm}
\end{answerbox}

\item Write the formulas for the \textbf{ATE} (eq.~[2.75]) and the \textbf{ATT} (eq.~[2.76]). In one sentence each, explain what population each one refers to.

\begin{answerbox}
\vspace{4cm}
\end{answerbox}
\end{enumerate}

\newpage
%=============================================================================
\section{Block 3: ATE vs.\ ATT and Selection Bias --- Numerical Exercise \textnormal{(13 min)}}
%=============================================================================

\textbf{Question 4} (Numerical exercise)

\medskip
Suppose we have a population with two types of individuals:

\begin{center}
\begin{tabular}{cccc}
\toprule
Type & Fraction of population & $y_i(0)$ & $y_i(1)$ \\
\midrule
A & 0.6 & 2 & 5 \\
B & 0.4 & 4 & 6 \\
\bottomrule
\end{tabular}
\end{center}

\begin{enumerate}[label=(\alph*)]
\item Compute the individual treatment effect $\tau$ for each type.

\begin{answerbox}
\vspace{2.5cm}
\end{answerbox}

\item Compute the \textbf{ATE}.

\begin{answerbox}
\vspace{3cm}
\end{answerbox}

\item Now suppose that \textbf{only Type~A} individuals get treated ($x = 1$ for Type~A, $x = 0$ for Type~B). Compute the \textbf{ATT}. Is it equal to the ATE? Why or why not?

\begin{answerbox}
\vspace{3.5cm}
\end{answerbox}

\item Compute the \textbf{selection bias} $E[y(0) \mid x{=}1] - E[y(0) \mid x{=}0]$. What is its sign? Interpret.

\begin{answerbox}
\vspace{3.5cm}
\end{answerbox}

\newpage
\item Verify the decomposition: compute both sides of
\[
  E[y \mid x{=}1] - E[y \mid x{=}0] = ATT + \text{Selection Bias}
\]

\begin{answerbox}
\vspace{4cm}
\end{answerbox}

\item Suppose instead that treatment were \textbf{randomly assigned} (each individual gets $x=1$ with probability 0.5, regardless of type). What would $E[\bar{y}_1 - \bar{y}_0]$ estimate?

\begin{answerbox}
\vspace{3.5cm}
\end{answerbox}
\end{enumerate}

%=============================================================================
\section{Block 4: The Selection Bias Decomposition and Random Assignment \textnormal{(13 min)}}
%=============================================================================

\textbf{Question 5} (Key derivation)

\medskip
We now prove the result we used numerically in Block~3. Let $\bar{y}_1$ be the average outcome among treated units and $\bar{y}_0$ among untreated.

\begin{enumerate}[label=(\alph*)]
\item Show that the \textbf{population} version of the comparison $\bar{y}_1 - \bar{y}_0$ decomposes as:
\[
  E[y \mid x{=}1] - E[y \mid x{=}0]
  = \underbrace{E\bigl[y(1) - y(0) \mid x{=}1\bigr]}_{ATT}
  + \underbrace{E\bigl[y(0) \mid x{=}1\bigr] - E\bigl[y(0) \mid x{=}0\bigr]}_{\text{Selection Bias}}
\]

\textcolor{hintcolor}{\emph{Hint: use the switching equation to write $E[y \mid x{=}1] = E[y(1) \mid x{=}1]$ and $E[y \mid x{=}0] = E[y(0) \mid x{=}0]$. Then add and subtract $E[y(0) \mid x{=}1]$.}}

\begin{answerbox}
\vspace{7cm}
\end{answerbox}

\newpage
\item Explain in words what $E[y(0) \mid x{=}1] - E[y(0) \mid x{=}0]$ means. Use the education example (Ex.~1.4) and the job training context to give one example of \textbf{positive} and one of \textbf{negative} selection bias.

\begin{answerbox}
\vspace{5cm}
\end{answerbox}

\item Now suppose treatment is \textbf{randomly assigned}. Before doing the math, let us understand \emph{why} randomization works.

\begin{bookexample}
\textbf{\color{blockcolor}Why does random assignment ``work''? --- The mathematical mechanism}

\smallskip
\textbf{Step 0: A key property of independence.} Recall that if $X \perp Y$, then
\[
  E[Y \mid X = x] = E[Y] \quad \text{for all } x.
\]
Conditioning on an independent variable \textbf{does nothing}: the conditional expectation equals the unconditional one. \emph{This single fact is the entire engine.}

\smallskip
\textbf{Step 1: Why does randomization create independence?} Before the experiment, each unit~$i$ already has \emph{fixed} potential outcomes $(y_i(0), y_i(1))$, determined by that person's characteristics (ability, motivation, health, \dots). These exist regardless of what we do.

When we randomly assign treatment, $x_i$ is determined by a coin flip that \textbf{knows nothing} about unit~$i$. Formally:
\[
  P(x_i = 1 \mid y_i(0),\, y_i(1)) = P(x_i = 1) = p \quad \text{for all } i
\]
A person with $y_i(0) = 100$ has the \emph{same} probability $p$ of being treated as a person with $y_i(0) = 2$. The coin does not look at the potential outcomes. Therefore, \textbf{by construction}:
\[
  \boxed{x \perp (y(0),\; y(1))}
\]

\smallskip
\textbf{Step 2: Apply the independence property.} Since $x \perp y(0)$:
\[
  E[y(0) \mid x{=}1] = E[y(0)] \quad\text{and}\quad E[y(0) \mid x{=}0] = E[y(0)]
\]
The treated group and the control group have \textbf{the same average baseline} --- not because they are identical person-by-person, but because the \emph{sorting mechanism} (the coin) is unrelated to individual characteristics. Likewise, $x \perp y(1)$ gives $E[y(1) \mid x{=}1] = E[y(1)]$.

\smallskip
\textbf{Step 3: Why does this fail without randomization?} Without randomization, people choose (or are selected for) treatment based on their characteristics --- the same characteristics that determine $(y_i(0), y_i(1))$. So $P(x_i = 1 \mid y_i(0), y_i(1)) \neq P(x_i = 1)$: the probability of treatment \emph{depends on} potential outcomes. For example, if high-ability people are more likely to go to college, then $E[y(0) \mid x{=}1] > E[y(0) \mid x{=}0]$. Conditioning on $x = 1$ selects a non-representative subgroup, and the independence property fails.
\end{bookexample}

Now use the property $x \perp (y(0), y(1))$ to show three things:
\begin{enumerate}[label=(\roman*)]
  \item Selection bias $= 0$.
  \item $ATT = ATE$.
  \item $\bar{y}_1 - \bar{y}_0$ is unbiased for the ATE.
\end{enumerate}

\begin{answerbox}
\vspace{7cm}
\end{answerbox}

\newpage
\item (Problem~15 from the book) The \textbf{sample average treatment effect} estimator is $\hat{\tau}_{ate} = n^{-1}\sum_{i=1}^{n}[y_i(1) - y_i(0)]$. Show that \emph{if} we could observe both potential outcomes for everyone, $\hat{\tau}_{ate}$ would be unbiased for $\tau_{ate} = E[y(1) - y(0)]$.

\begin{answerbox}
\vspace{6cm}
\end{answerbox}

\item Explain briefly why $\bar{y}_1$ and $\bar{y}(1) \equiv n^{-1}\sum_{i=1}^{n} y_i(1)$ are \emph{not} the same thing.

\begin{answerbox}
\vspace{4cm}
\end{answerbox}
\end{enumerate}

%=============================================================================
\section{Block 5: Application --- Job Training Program \textnormal{(6 min)}}
%=============================================================================

\textbf{Question 6} (From Example~2.14 in the book)

\begin{bookexample}
\textbf{\color{blockcolor}Example 2.14 --- Evaluating a Job Training Program (JTRAIN2)}

\smallskip
The data in JTRAIN2 are from the National Supported Work demonstration, where men were \textbf{randomly assigned} to receive job training (9--18 months) or serve as controls. The outcome is \texttt{re78}: real earnings in 1978 (thousands of 1982 dollars). The treatment variable is \texttt{train} ($= 1$ if trained).

Of 445 men, 185 received training and 260 did not. The simple regression $\widehat{re78} = \hat{\beta}_0 + \hat{\beta}_1 \cdot train$ gives: $\hat{\beta}_0 = 4.555$ (control group average, in thousands), $\hat{\beta}_1 = 1.794$ (treated earned \$1{,}794 more on average). The $t$-statistic is 1.79, and $R^2 \approx 0.018$.

The small $R^2$ means training explains only 1.8\% of earnings variation. But $R^2$ measures \textbf{explanatory power}, not causal significance --- under random assignment, $\hat{\beta}_1$ is still unbiased for the ATE.
\end{bookexample}

\begin{enumerate}[label=(\alph*)]
\item Express $\bar{y}_1$, $\bar{y}_0$, and $\hat{\beta}_1$ in terms of each other. Compute $\bar{y}_1$.

\begin{answerbox}
\vspace{3.5cm}
\end{answerbox}

\item Can we interpret $\hat{\beta}_1 = 1.794$ causally? Why? Express the effect as a percentage of the control mean.

\begin{answerbox}
\vspace{3.5cm}
\end{answerbox}

\item The $R^2$ is 0.018. A student says: ``The training had no effect because $R^2$ is tiny.'' Is this correct?

\begin{answerbox}
\vspace{3.5cm}
\end{answerbox}

\item If workers had \textbf{volunteered} instead of being randomly assigned, would you still trust the estimate? Why?

\begin{answerbox}
\vspace{3.5cm}
\end{answerbox}
\end{enumerate}

\end{document}
